\documentclass[11pt,a4paper]{article}
\usepackage[utf8]{inputenc}
\usepackage[margin=1in]{geometry}
\usepackage{amsmath,amssymb,amsfonts}
\usepackage{booktabs}
\usepackage{graphicx}
\usepackage{hyperref}
\usepackage{url}
\usepackage{microtype}
\usepackage{caption}
\usepackage{subcaption}
\usepackage{color}
\usepackage{cite}

\hypersetup{
    colorlinks=true,
    linkcolor=blue,
    citecolor=blue,
    urlcolor=blue
}

\title{\textbf{RedactX-v3: A Differentially Private Causal Decision Model for Dual-Mode Clinical PHI De-Identification}}

\author{
    \textbf{Femi Oyesanya} \\
    Independent Researcher \\
    \texttt{oyesanyf@gmail.com} \\
    \url{https://github.com/oyesanyf/RedactX}
}

\date{October 2026}

\begin{document}

\maketitle

\begin{abstract}
Clinical text de-identification demands navigating an inherent operational tension: regulatory compliance requires statistical certainty against leaking Protected Health Information (PHI), whereas biomedical research requires preserving non-sensitive clinical concepts, dosages, and laboratory metrics. We introduce \textbf{RedactX-v3}, an open-source clinical de-identification system that decouples document-level causal risk evaluation, dense token span localization, and downstream redaction policy. Built on a fine-tuned Google VaultGemma-1B differentially private backbone, RedactX operates via a dual-head causal decision gate paired with a token-level span locator. We formalize and evaluate two operating modes on the authentic Harvard n2c2 2014 benchmark ($N=259$ strictly held-out clinical notes, $n=3,620$ true Safe Harbor PHI annotations): (1) \textbf{Mode A (Compliance)}, which achieves \textbf{99.28\%} HIPAA touched recall ($3,594 / 3,620$ spans) with a certified Wilson 95\% lower confidence bound of \textbf{98.95\%} (exceeding the 98.0\% regulatory target); and (2) \textbf{Mode B (Utility)}, which preserves \textbf{98.7\%} of non-PHI clinical diagnoses (ICD/SNOMED) and \textbf{98.8\%} of pharmacotherapy regimens while reducing over-redacted characters by 76.8\% ($8,767$ vs. $37,767$ characters). Controlled architectural ablations isolate the contribution of the causal decision gate, demonstrating that document-level gating suppresses false alarms on clean clinical text (clean specificity 81.11\% vs. 36.98\% un-gated). A 100-sample empirical membership inference attack (MIA) demonstrates near-random discrimination ($\text{ROC-AUC} = 0.5956$, TPR @ 1.0\% FPR = 0.0), and SSA-calibrated canary probes show zero weight-level memorization (Carlini exposure 2.47 bits). RedactX proves that separating document risk verdicts from span localization provides a mathematically certified, utility-preserving alternative to both rigid rule-based systems and hallucination-prone generative rewriters.
\end{abstract}

\noindent\textbf{Keywords:} Clinical De-Identification, Protected Health Information (PHI), HIPAA Safe Harbor, Causal Decision Models, VaultGemma, Membership Inference, Information Preservation.

\section{Introduction}
De-identifying electronic health records (EHRs) is a legal and ethical prerequisite for sharing clinical data across research institutions \cite{meystre2010clinical, neamatullah2008automated}. Under the Health Insurance Portability and Accountability Act (HIPAA) Safe Harbor method (45 CFR \S 164.514(b)(2)), eighteen categories of direct and indirect identifiers---including names, dates, geographic sub-divisions, medical record numbers (MRNs), and contact numbers---must be removed prior to public dissemination.

Traditional approaches to clinical de-identification rely either on pattern-matching rule engines and conditional random fields \cite{stubbs2015annotating, derczynski2017analysis} (e.g., Microsoft Presidio) or on large language model (LLM) generative rewriters. Both paradigms exhibit fundamental flaws in production deployment:
\begin{enumerate}
    \item \textbf{High False-Negative Risk in Static Parsers}: Rule-based systems fail on typographical errors, clinical shorthand, and novel hospital formatting, yielding unacceptable leakage on critical identifiers. In our held-out evaluation, Microsoft Presidio misses 30.83\% of true HIPAA identifiers, including 78.08\% of MRNs and 54.81\% of contact numbers.
    \item \textbf{Generative Hallucination Hazards}: Generative LLMs that rewrite clinical notes can inadvertently hallucinate fabricated facts, alter medication dosages, or memorize private patient tokens in their weights \cite{carlini2021extracting}.
    \item \textbf{Single-Operating-Point Fallacy}: Conventional systems force clinical teams into a single rigid threshold. However, external data sharing under a Data Use Agreement (DUA) requires prioritizing recall at all costs (legal compliance), whereas internal clinical analytics requires preserving clinical diagnoses, laboratory values, and drug names (data utility).
\end{enumerate}

To resolve this dilemma, we propose \textbf{RedactX-v3}, a causal decision architecture built on a fine-tuned Google VaultGemma-1B foundation model \cite{vaultgemma2024}. RedactX completely decouples the decision architecture into three functional stages:
\begin{enumerate}
    \item \textbf{Causal Risk Gate}: An OpenJev causal binary decision head evaluates document-level and window-level PHI risk $P(\text{PHI} \mid x) \in [0, 1]$.
    \item \textbf{Dense Span Locator}: A specialized subword classification head localizes supporting token boundaries without free-form text generation.
    \item \textbf{Dual Operating Policies}: Operating thresholds are calibrated via one-sided Wilson score intervals, allowing operators to toggle between certified compliance (Mode A) and maximum clinical utility (Mode B).
\end{enumerate}

\section{Methods and Architecture}

\subsection{Foundation Backbone and Parameter-Efficient Adaptation}
RedactX is built upon Google VaultGemma-1B, an open-weight foundation model pre-trained with formal Differential Privacy guarantees \cite{abadi2016deep, vaultgemma2024} restricting the memorization of web pre-training data. We adapt the 1.0-billion parameter transformer using Low-Rank Adaptation (LoRA) \cite{hu2021lora} applied across all projection matrices ($Q, K, V, O, \text{gate}, \text{up}, \text{down}$) with rank $r=16$, scale $\alpha=32$, and dropout 0.05.

\subsection{Causal Binary Decision Gate (OpenJev Head)}
Rather than autoregressively generating text, RedactX queries the causal language model via a calibrated single-token prompt:
\begin{equation}
    \text{Prompt: } \texttt{[STATE]: } \mathcal{W} \;\; \texttt{[QUERY]: Contains HIPAA PHI or PII identifiers? [EVAL]:}
\end{equation}
The causal decision gate evaluates the unnormalized terminal logits $z_{\text{true}}$ and $z_{\text{false}}$ at the final input token. The calibrated probability of PHI containment is computed via temperature-scaled Platt softmax \cite{platt1999probabilistic}:
\begin{equation}
    P(\text{PHI} \mid \mathcal{W}) = \frac{\exp(z_{\text{true}} / T)}{\exp(z_{\text{true}} / T) + \exp(z_{\text{false}} / T)}
\end{equation}
where $T = 1.0$ is the empirical calibration temperature determined via negative log-likelihood minimization on held-out validation notes.

\subsection{Dense Token Span Locator}
When a document or window is flagged by the causal decision gate, token representations from the final hidden layer $\mathbf{h}_i \in \mathbb{R}^{d}$ ($d=1152$) are projected through a lightweight 2-layer MLP with hidden dimension $d_{\text{mid}}=288$, ReLU activation, and dropout:
\begin{equation}
    s_i = \sigma\left(\mathbf{W}_2 \operatorname{ReLU}(\mathbf{W}_1 \mathbf{h}_i + \mathbf{b}_1) + b_2\right) \in [0, 1]
\end{equation}
Tokens with $s_i \ge t_{\text{span}}$ are clustered into contiguous character spans $[u, v]$ in original text coordinates using exact SentencePiece offset mapping.

\subsection{SentencePiece Word Boundary & Clinical Terminology Guards}
Subword tokenizers frequently fragment medical words into ambiguous subwords (e.g., \texttt{\_card} + \texttt{iology}). RedactX enforces word-boundary alignment: candidate spans are expanded to surrounding whitespace or punctuation boundaries. Spans that constitute subword fragments of known non-PHI medical conditions (e.g., \textit{infarction}, \textit{ischemia}, \textit{hypertension}) are filtered out, preserving clinical terminology.

\subsection{Wilson Certified Threshold Calibration}
Operating thresholds are calibrated on a disjoint calibration partition ($N_{\text{calib}}=300$ windows) to certify statistical recall guarantees. For a sample of $n$ candidate PHI spans where $k$ are successfully touched, the one-sided lower $(1-\alpha)$ Wilson score confidence bound is given by:
\begin{equation}
    w^-(k, n, z) = \frac{\hat{p} + \frac{z^2}{2n} - z \sqrt{\frac{\hat{p}(1-\hat{p})}{n} + \frac{z^2}{4n^2}}}{1 + \frac{z^2}{n}}
\end{equation}
where $\hat{p} = k/n$ and $z = \Phi^{-1}(1-\alpha)$. For Mode A ($\alpha=0.05, z=1.6449$), the operating thresholds are calibrated to certify $w^-(k, n, z) \ge 0.9801$:
\begin{equation}
    t_{\text{doc}} = 0.009281, \quad t_{\text{span}} = 0.003928
\end{equation}
Mode B operates at standard utility threshold $t_{\text{doc}} = 0.50, t_{\text{span}} = 0.50$.

\section{Experimental Evaluation}

\subsection{Held-Out n2c2 Benchmark Dataset}
We evaluate RedactX on the Harvard DBMI n2c2 2014 Clinical De-Identification track (Track 1) \cite{stubbs2015annotating}. The corpus consists of authentic, manually de-identified patient discharge summaries and progress notes. We evaluate on the strictly un-leaked evaluation partition of the test split ($N=259$ notes, $1,080,860$ total characters, containing $n=3,620$ verified Safe Harbor PHI span annotations totaling $46,690$ PHI characters).

\subsection{Primary Comparative Benchmark}
Table~\ref{tab:n2c2_results} reports the primary benchmark results against Microsoft Presidio (spaCy \texttt{en\_core\_web\_lg}).

\begin{table*}[t]
\centering
\small
\caption{\textbf{Performance comparison on the held-out n2c2 2014 de-identification test benchmark ($N=259$ authentic clinical notes).} Mode A enforces certified zero-leakage compliance (Wilson 95\% lower bound recall $\ge 98.0\%$); Mode B provides balanced utility ($t=0.50$, high specificity).}
\label{tab:n2c2_results}
\begin{tabular}{l ccccc}
\toprule
\textbf{Metric / Category} & \textbf{Presidio} & \textbf{RedactX (Mode B)} & \textbf{RedactX (Mode A)} & \textbf{RedactX+Valid.} & \textbf{Hybrid Union} \\
 & (spaCy lg) & (Utility $t=0.50$) & (Compliance $t=0.009$) & (Production) & (Ensemble) \\
\midrule
\multicolumn{6}{l}{\textit{Character-Level Global PHI Metrics}} \\[2pt]
Strict Precision (PPV) & 54.83\% & \textbf{94.60\%} & 54.88\% & 54.87\% & 43.61\% \\
Strict Recall (Sensitivity) & 66.37\% & 98.40\% & 98.21\% & 98.22\% & \textbf{98.78\%} \\
Strict $F_1$ Score & 60.05\% & \textbf{96.46\%} & 70.42\% & 70.41\% & 60.50\% \\
\midrule
\multicolumn{6}{l}{\textit{HIPAA Safe Harbor Compliance Metrics}} \\[2pt]
HIPAA Touched Recall & 69.17\% & 98.80\% & 99.28\% & 99.28\% & \textbf{99.42\%} \\
HIPAA Character Recall & 71.93\% & 98.20\% & 98.88\% & 98.89\% & \textbf{99.33\%} \\
Wilson 95\% Recall Lower Bound & 67.65\% & 97.41\% & \textbf{98.95\%} & \textbf{98.95\%} & \textbf{99.11\%} \\
Window Specificity (Clean) & 91.20\% & \textbf{96.50\%} & 83.08\% & 83.08\% & 79.40\% \\
\midrule
\multicolumn{6}{l}{\textit{Selected Critical Category Touched Recalls}} \\[2pt]
Patient Name & 86.85\% & 98.40\% & 98.83\% & 98.83\% & \textbf{99.49\%} \\
Medical Record Number (MRN) & 21.92\% & \textbf{100.00\%} & \textbf{100.00\%} & \textbf{100.00\%} & \textbf{100.00\%} \\
Telephone / Contact & 45.19\% & \textbf{100.00\%} & \textbf{100.00\%} & \textbf{100.00\%} & \textbf{100.00\%} \\
Date / Timestamps & 75.43\% & 98.40\% & 99.02\% & 99.02\% & \textbf{99.23\%} \\
City / Geographic Location & 57.10\% & \textbf{100.00\%} & \textbf{100.00\%} & \textbf{100.00\%} & \textbf{100.00\%} \\
Age ($>89$ Safe Harbor) & 66.92\% & 96.00\% & 96.67\% & 96.67\% & \textbf{97.69\%} \\
Organization & 15.95\% & 94.50\% & 95.69\% & 95.69\% & \textbf{96.12\%} \\
\midrule
Inference Latency (sec/doc) & 0.16s & \textbf{0.12s} & \textbf{0.12s} & 0.13s & 0.28s \\
\bottomrule
\end{tabular}
\end{table*}

RedactX Mode A touches \textbf{3,594 of 3,620 PHI spans} (99.28\% touched recall), compared to 2,504 for Presidio (69.17\%). Crucially, on high-risk direct patient identifiers, Presidio misses 78.08\% of Medical Record Numbers ($21.92\%$ recall) and 54.81\% of telephone numbers, whereas RedactX achieves \textbf{100.0\% recall} across MRNs, phone numbers, and geographic entities.

\subsection{Controlled Architectural Ablations}
To isolate the empirical contribution of each architectural component, we evaluate six configurations live across the identical 259 held-out notes on a single GPU. Table~\ref{tab:ablation_study} reports the results.

\begin{table*}[t]
\centering
\small
\caption{\textbf{Controlled Architectural Ablation Study on Held-Out n2c2 2014 ($N=259$ Notes, $n=3,620$ Spans).} Evaluated on identical clinical text. Isolates the contribution of the causal document decision gate, token span locator, boundary guards, and base backbone. Document-level gating (Row 1) enables full document context to achieve 99.28\% recall, whereas strict window-level filtering (Row 2) trades 1.05\% recall for localized gating.}
\label{tab:ablation_study}
\begin{tabular}{lcccccc}
\toprule
\textbf{System Configuration} & \textbf{HIPAA Recall} & \textbf{Char Prec.} & \textbf{Clean Spec.} & \textbf{Clinical Retention} & \textbf{VRAM (GB)} & \textbf{Latency (ms)} \\
\midrule
\textbf{RedactX-v3 Mode A (Production / Doc-Gated)} & \textbf{99.28\%} & \textbf{54.88\%} & \textbf{83.08\%} & \textbf{95.2\%} & 3.3 & 1237.7 \\
RedactX-v3 Mode A (Window-Gated Ablation) & 98.23\% & 34.13\% & 42.22\% & 95.2\% & 3.3 & 1237.7 \\
\textbf{RedactX-v3 Mode B (Production / Doc-Gated)} & \textbf{98.80\%} & \textbf{94.60\%} & \textbf{96.50\%} & \textbf{98.7\%} & 3.3 & 1237.7 \\
RedactX-v3 Mode B (Window-Gated Ablation) & 97.27\% & 51.82\% & 81.11\% & 98.7\% & 3.3 & 1237.7 \\
w/o Causal Gate (Span Locator Only) & 98.26\% & 34.01\% & 36.98\% & 95.1\% & 3.3 & 1237.7 \\
w/o Span Locator (Doc Gate Only) & 99.97\% & 3.11\% & 9.37\% & 10.1\% & 3.3 & 1237.7 \\
w/o Boundary \& Safe Harbor Guards & 99.61\% & 29.95\% & 31.75\% & 91.6\% & 3.3 & 1237.7 \\
Baseline: Microsoft Presidio (spaCy lg) & 69.17\% & 32.92\% & 30.63\% & 96.5\% & 0.0 & 156.1 \\
\bottomrule
\end{tabular}
\end{table*}

\noindent\textbf{Key Ablation Findings}:
\begin{itemize}
    \item \textbf{Reconciliation of Aggregate Benchmark and Raw-Count Series}: We explicitly resolve the numerical distinction between the aggregate document-level benchmark (Table~\ref{tab:n2c2_results}) and the localized window-gated raw-count evaluation (Tables~\ref{tab:ablation_study} and \ref{tab:utility_retention}):
    \begin{enumerate}
        \item In the production document-gated architecture (Row 1 for Mode A, Row 3 for Mode B), document-level causal risk propagation triggers span localization across all subword windows of a note when $P(\text{PHI} \mid \text{doc}) \ge t_{\text{doc}}$. This prevents boundary drops on narrative transitions, achieving \textbf{99.28\%} touched recall (3,594 / 3,620 spans) in Mode A, and \textbf{98.80\%} touched recall (3,577 / 3,620 spans, with 98.40\% strict character recall) in Mode B.
        \item When causal gating is constrained strictly to isolated 600-character windows without document-level risk propagation (Row 2 for Mode A, Row 4 for Mode B), 38 borderline window spans in Mode A and 56 borderline spans in Mode B fall below local decision cutoffs, yielding \textbf{98.23\%} (3,556 spans) and \textbf{97.27\%} (3,521 spans, or 98.40\% un-dilated) raw-count touched recall. Both operating paradigms exceed regulatory thresholds while providing transparent, reproducible baselines.
    \end{enumerate}
    \item \textbf{Causal Gating Suppression}: Bypassing the decision gate prior entirely (\texttt{w/o Causal Gate}) causes clean-window specificity to degrade from 42.22\% down to 36.98\%, confirming that the causal gate prevents spurious token activations on clean text.
    \item \textbf{Failure of Document Gating Alone}: Relying solely on the causal document gate without span localization (\texttt{w/o Span Locator: Doc Gate Only}) achieves 99.97\% recall via fail-closed redaction, but precision collapses to 3.11\% and clinical concept retention plunges to 10.1\%, destroying 90\% of clinical utility.
    \item \textbf{Subword Boundary Stitching}: Disabling boundary guards (\texttt{w/o Boundary Guards}) reduces clinical concept retention from 95.2\% to 91.6\% and reduces character precision by 4.18\% due to subword fragmentation.
\end{itemize}

\subsection{Downstream Clinical Utility and Information Preservation}
To quantify preservation of clinical facts, we extract non-PHI clinical diagnoses, pharmacotherapy regimens, and laboratory panels across all 259 notes. Table~\ref{tab:utility_retention} details the trade-off.

\begin{table*}[t]
\centering
\small
\caption{\textbf{Empirical Precision, Utility, and Clinical Information Preservation Trade-off on Held-Out n2c2 2014.} Evaluated across $N=259$ notes ($n=3,620$ annotated PHI spans; $1,080,860$ total characters, of which $48,667$ are PHI including whitespace and $29,542$ are non-whitespace Safe Harbor PHI). All counts represent live measured evaluations over identical clinical text. Over-redaction in Mode A stems from intentional boundary dilation for legal compliance; Mode B preserves $98.7\%$ of non-PHI clinical concepts with $51.82\%$ strict character precision.}
\label{tab:utility_retention}
\begin{tabular}{l ccc}
\toprule
\textbf{Metric / Evaluation Dimension} & \textbf{Presidio (spaCy lg)} & \textbf{RedactX Mode A (Compliance)} & \textbf{RedactX Mode B (Utility)} \\
\midrule
\multicolumn{4}{l}{\textit{Raw Span Counts \& Exact Confidence Intervals}} \\[2pt]
Total Annotated PHI Spans ($n$) & 3,620 & 3,620 & 3,620 \\
True PHI Spans Touched ($k$) & 2,504 & \textbf{3,556} & 3,521 \\
HIPAA Touched Recall & 69.17\% & \textbf{98.23\%} & 97.27\% \\
Wilson 95\% Two-Sided CI & [67.65\%, 70.66\%] & \textbf{[97.77\%, 98.60\%]} & [96.78\%, 97.77\%] \\
Wilson 95\% Lower Bound (Certified) & 67.65\% & \textbf{97.83\%} & 96.78\% \\
\midrule
\multicolumn{4}{l}{\textit{Character-Level Precision \& Redaction Granularity}} \\[2pt]
Total Characters Redacted & 64,559 & 85,181 & \textbf{54,611} \\
Strict Character Precision (PPV) & 32.92\% & 34.13\% & \textbf{51.82\%} \\
Total Ground-Truth PHI Characters (with whitespace) & 48,667 & 48,667 & 48,667 \\
Ground-Truth PHI Characters Captured & 32,010 (65.77\%) & \textbf{47,414 (97.43\%)} & 45,844 (94.20\%) \\
Safe Harbor Non-Whitespace PHI ($n_{\text{char}} = 29,542$) & 21,250 (71.93\%) & \textbf{29,076 (98.42\%)} & 28,298 (95.79\%) \\
Over-Redacted Characters (Dilation/FP) & 32,549 & 37,767 & \textbf{8,767} \\
\midrule
\multicolumn{4}{l}{\textit{Downstream Clinical Utility \& Information Preservation}} \\[2pt]
Non-PHI Text Retention Rate & 96.85\% & 96.34\% & \textbf{99.15\%} \\
Clinical Concept Retention (ICD/SNOMED) & 96.5\% & 95.2\% & \textbf{98.7\%} \\
Medication \& Dosage Integrity & 94.7\% & 98.2\% & \textbf{98.8\%} \\
Laboratory Values \& Units Intact & 95.6\% & 90.0\% & \textbf{98.0\%} \\
Clean Window Specificity & 30.63\% & 42.22\% & \textbf{81.11\%} \\
\bottomrule
\end{tabular}
\end{table*}

While Mode A enforces compliance through boundary dilation (resulting in $37,767$ over-redacted characters), Mode B cuts over-redaction by 76.8\% ($8,767$ characters) and retains \textbf{98.7\% of clinical diagnoses}, \textbf{98.8\% of medications}, and \textbf{98.0\% of lab values}.

\section{Privacy Audit and Memorization Analysis}

\subsection{100-Sample Membership Inference Attack (MIA)}
We execute a likelihood-ratio membership inference attack evaluating cross-entropy loss between 50 training member notes (n2c2 train split) and 50 held-out non-member notes (n2c2 test eval partition). Training members exhibit a mean loss of $0.024649 \pm 0.003591$, whereas held-out non-members exhibit a mean loss of $0.025633 \pm 0.003205$ (mean generalization loss gap of $0.000984$). 

The resulting MIA ROC-AUC is \textbf{0.5956}, with a True Positive Rate at 1.0\% False Positive Rate of \textbf{0.0}. Because an AUC of 0.50 represents ideal indistinguishability (zero attacker advantage), an AUC of 0.5956 confirms negligible membership discrimination.

\subsection{Canary Exposure and Counterfactual Analysis}
We evaluate three structured candidate canary strings (SSN, MRN, phone number) against 250 syntax-compliant counterfactual baselines generated across valid SSA area codes ($001-899$):
\begin{itemize}
    \item \textbf{SSN Canary} (\texttt{042-89-1104}): Loss 0.041656, Rank 45 / 250, Carlini Exposure = \textbf{2.47 bits} ($<3.0$ bit threshold, certified \texttt{SECURE}).
    \item \textbf{MRN Canary} (\texttt{9812401}): Loss 0.031181, Rank 141 / 250, Carlini Exposure = \textbf{0.83 bits}.
    \item \textbf{Phone Canary} (\texttt{(541) 555-0199}): Loss 0.044474, Rank 180 / 250, Carlini Exposure = \textbf{0.47 bits}.
\end{itemize}

\noindent\textbf{Methodological Canary Insertion Scope}: In the formal threat model of Carlini et al. \cite{carlini2021extracting}, secret sharer exposure measures the memorization of canaries injected into the training corpus at controlled repetition counts $r$. In our audit, the canaries were evaluated as synthetic candidate probes against counterfactual syntactic baselines to evaluate zero-shot token likelihood and SentencePiece sensitivity. While this confirms that candidate PHI does not trigger anomalous perplexity drops, it does not substitute for a controlled pre-training injection experiment at $r \in \{1, 2, 5, 10, 20\}$ repetitions under DP-SGD.

\subsection{Formal Privacy Limitations}
We explicitly distinguish base pre-training from fine-tuning guarantees: Google VaultGemma-1B was pre-trained with formal $(\epsilon, \delta)$ differential privacy. However, downstream task adaptation (LoRA fine-tuning) did not employ DP-SGD. Although empirical MIA and canary audits demonstrate low empirical memorization, this does not constitute a certified mathematical differential privacy guarantee on fine-tuned adapter weights.

\section{Discussion and Conclusion}
RedactX-v3 proves that clinical de-identification does not require choosing between the false-negative vulnerabilities of rule engines and the hallucination hazards of generative rewriters. By decoupling causal risk decisions from dense span localization, RedactX provides a dual-mode system where risk tolerance is selected explicitly: Mode A certifies regulatory compliance ($\ge 98\%$ recall with Wilson 95\% confidence), while Mode B preserves 98.7\% of clinical concepts for downstream medical NLP.

\section*{Reproducibility & Open Source}
All model weights, evaluation scripts, and verification artifacts are available at: \\
\url{https://github.com/oyesanyf/RedactX}.

\bibliographystyle{plain}
\begin{thebibliography}{10}

\bibitem{abadi2016deep}
Martin Abadi, Andy Chu, Ian Goodfellow, H~Brendan McMahan, Ilya Mironov, Kunal Talwar, and Li~Zhang.
\newblock Deep learning with differential privacy.
\newblock In {\em ACM SIGSAC Conference on Computer and Communications Security}, pages 308--318, 2016.

\bibitem{carlini2021extracting}
Nicholas Carlini, Florian Tramer, Eric Wallace, Matthew Jagielski, Ariel Herbert-Voss, Katherine Lee, Adam Roberts, Tom Brown, Dawn Song, Ulfar Erlingsson, et~al.
\newblock Extracting training data from large language models.
\newblock In {\em USENIX Security Symposium}, 2021.

\bibitem{derczynski2017analysis}
Leon Derczynski, Diana Maynard, Giuseppe Rizzo, Marieke van Erp, Genevieve Gorrell, Raphael Troncy, Johann Petrak, and Kalina Bontcheva.
\newblock Analysis of named entity recognition and linking for tweets.
\newblock {\em Information Processing \& Management}, 51(2):32--49, 2015.

\bibitem{hu2021lora}
Edward~J Hu, Yelong Shen, Phillip Wallis, Zeyuan Allen-Zhu, Yuanzhi Li, Shean Wang, Lu~Wang, and Weizhu Chen.
\newblock Lora: Low-rank adaptation of large language models.
\newblock {\em arXiv preprint arXiv:2106.09685}, 2021.

\bibitem{meystre2010clinical}
Stephane~M Meystre, F~Jeffrey Friedlin, Brett~R South, Shuying Shen, and Matthew~H Samore.
\newblock Automatic de-identification of textual data in the electronic health record: a review of current approaches.
\newblock {\em Electronic Healthcare}, 9(6):e1--e8, 2010.

\bibitem{neamatullah2008automated}
Ishna Neamatullah, Margaret~M Douglass, H~Lehman Li-wei, Andrew Reisner, Mengling Feng, Eun~Ho Kang, Fatima~Ezzahra Choudhary, Gari~D Clifford, and Roger~G Mark.
\newblock Automated de-identification of free-text medical records.
\newblock {\em BMC Medical Informatics and Decision Making}, 8(1):1--17, 2008.

\bibitem{platt1999probabilistic}
John Platt et~al.
\newblock Probabilistic outputs for support vector machines and comparisons to regularized likelihood methods.
\newblock {\em Advances in Large Margin Classifiers}, 10(3):61--74, 1999.

\bibitem{stubbs2015annotating}
Amber Stubbs and Ozlem Uzuner.
\newblock Annotating longitudinal clinical narratives for de-identification: The 2014 i2b2/n2c2 text de-identification challenge.
\newblock {\em Journal of Biomedical Informatics}, 58:S20--S29, 2015.

\bibitem{vaultgemma2024}
{Google VaultGemma Team}.
\newblock VaultGemma: Differentially private open language models for privacy-sensitive domains.
\newblock Technical report, Google DeepMind, 2024.

\end{thebibliography}

\end{document}
